\documentclass[runningheads,a4paper]{llncs}

\usepackage{amsmath,amssymb}
\usepackage{graphicx}
\usepackage{booktabs}
\usepackage{multirow}
\usepackage{siunitx}
\usepackage{url}
\usepackage{hyperref}
\usepackage{xcolor}

\hypersetup{
    pdftitle={Energy-Aware Priority Scheduling in Smart Sensor Networks Using Policy-Guided Deep Reinforcement Learning},
    pdfauthor={Anand Emmanuel Raju B., Radhika N., Radhika G.},
    pdfsubject={ICTSES 2026 Camera-Ready Manuscript},
    pdfkeywords={Smart Energy Systems, Wireless Sensor Networks, Reinforcement Learning, ns-3, QoS, Priority Scheduling}
}

\title{Energy-Aware Priority Scheduling in Smart Sensor Networks Using Policy-Guided Deep Reinforcement Learning}
\titlerunning{Energy-Aware Priority Scheduling in Smart Sensor Networks}

\author{Anand Emmanuel Raju B.\inst{1} \and Radhika N.\inst{1} \and Radhika G.\inst{1}}
\authorrunning{A. E. R. Billa et al.}

\institute{
Department of Computer Science and Engineering,\\
Amrita School of Computing, Coimbatore,\\
Amrita Vishwa Vidyapeetham, India\\
\email{cb.sc.p2cse25003@cb.students.amrita.edu, n\_radhika@cb.amrita.edu, g\_radhika@cb.amrita.edu}
}

\begin{document}

\maketitle

\begin{abstract}
Smart grid sensor networks face a fundamental challenge: reliably delivering critical protection alerts (fault notifications, breaker commands) to operators while maintaining continuous telemetry feeds (voltage, current, temperature). Strict Priority scheduling---always transmit alarms first---achieves superior alert latency but starves routine telemetry during disturbances, exactly when operators need full visibility. This work investigates whether deep reinforcement learning can balance competing priorities better than rule-based approaches. A Double-DQN agent was trained using ns-3 simulation coupled with PyTorch and evaluated across 20 matched trials. Results quantify a clear trade-off: Strict Priority delivers critical packets 5.2 times faster (0.0274\,s vs. 0.1431\,s), whereas the learned scheduler reduces telemetry latency by 30.5\% (0.4287\,s vs. 0.6171\,s) at the cost of 2.56\% lower critical-packet delivery. For operators managing sustained grid disturbances, this fairness improvement supports situational awareness when it matters most.

\keywords{Smart Energy Systems \and Wireless Sensor Networks \and Reinforcement Learning \and ns-3 \and QoS \and Priority Scheduling}
\end{abstract}

\section{Introduction}
Power grids at the distribution level depend heavily on wireless sensor networks to do three things: detect equipment faults, monitor grid health in real time, and enable remote control of breakers. In practice, every sensor gateway has to simultaneously handle two very different kinds of traffic. One is urgent: protection alerts that must get through in milliseconds. The other is routine but important: telemetry streams (voltage, current, temperature) that operators integrate into their state-estimation algorithms.

Today's deployed systems handle this with Strict Priority scheduling. The rule is straightforward: critical packets always go first; routine telemetry is sent only when the critical queue is empty. It works respectably for critical alerts---operators do receive their alarms. But during actual grid faults, when protection alerts come in intensely and sustain for seconds, the routine telemetry queue starves completely. Operators watching the control room see nothing but alarms and lose visibility into basic grid conditions. This creates a critical gap: operators lose situational awareness exactly when they need it most during emergencies.

Can machine learning do better? The appeal is obvious: instead of a rigid rule, let a gateway learn an adaptive policy that carefully balances the two traffic objectives. A learning agent might recognize that a newly arrived critical alert doesn't require immediate transmission if its deadline is still 500 milliseconds away---that window could be used to clear aging telemetry packets before they expire. But this is hard to get right. Naive learning often leads to policy collapse: the agent learns to ignore all critical traffic and just race through routine packets to maximize immediate rewards.

Our work addresses this challenge directly. We designed a Double-DQN agent with deliberate reward shaping to avoid collapse, trained it in ns-3 simulation coupled with PyTorch, and ran 20 careful matched trials against Strict Priority and Random baselines. The key finding is honest and practical: our learned scheduler does not outperform Strict Priority on alert latency. Critical packets take 0.1431 seconds under our agent versus 0.0274 seconds under Strict Priority. But there's a genuine trade-off. Our agent cuts telemetry latency by 30.5\% (from 0.6171s to 0.4287s). For operators managing sustained grid faults, this responsiveness improvement carries practical value, while Strict Priority remains faster for isolated critical alerts. 

The core contributions include:
\begin{itemize}
    \item A fully reproducible simulation pipeline coupling ns-3.40 with PyTorch, enabling direct policy comparison under controlled conditions. The infrastructure is open-source and extensible.
    \item Soft reward guidance that prevents policy collapse while maintaining dual-objective flexibility. Guidance intensity varies with packet age, relaxing when deadlines have ample margin and tightening as expiration approaches.
    \item Empirical evidence that the learned policy generalizes reasonably across three network scales (50, 100, 200 nodes), though performance remains scenario-dependent.
    \item Quantified trade-offs: 2.56\% reduction in critical delivery (98.34\% to 95.78\%) exchanged for 30.5\% telemetry latency improvement. This fairness ratio favors operators requiring sustained awareness during grid disturbances.
\end{itemize}

\section{Related Work}
Reinforcement learning has seen application to wireless communications over the past decade. The motivation is straightforward: rather than manually configure network protocols, learning agents can adapt to varying conditions. Comprehensive surveys \cite{luong2019drlsurvey,frikha2021wrlsurvey} cover work on spectrum allocation, routing, MAC design, and packet scheduling. However, both surveys note a consistent challenge: RL methods require large training datasets, provide no formal safety assurances, and frequently fail when deployed systems differ significantly from training conditions.

Within sensor networks, several implementations and analyses bear on the problem. Leong \cite{leong2020automatica} demonstrated that neural policies adapt to measured packet arrival distributions and fading patterns that fixed heuristics cannot track. Sharma and Chauhan \cite{sharma2020wsn} explored distributed RL for simultaneous coverage and lifetime goals. Safety aspects appear in Mishra and Liang's Q-learning sleep optimization \cite{mishra2023sleep}, Nagib's deployment acceleration techniques \cite{nagib2023safe}, and Kim's work on real-time deadline satisfaction \cite{kim2024wtsn}.

A gap remains in the literature: direct comparison between learned schedulers and the rule-based systems they aim to replace. Most evaluations contrast learning against random assignment or no control, not against production heuristics like Strict Priority. We address this gap by benchmarking the learned policy against the industry standard scheduling approach, quantifying the operational trade-off and identifying where learning adds value.

\section{System Model and Problem Formulation}

\subsection{Network and Traffic Model}
We model an edge sensing cluster representative of a smart grid distribution substation. The network comprises 100 wireless sensor nodes with uniform random distribution across 150\,m $\times$ 150\,m, communicating with a central sink gateway. The physical and MAC layers employ ns-3.40 implementation of IEEE 802.15.4 in the 2.4\,GHz band.

Sensor nodes generate two distinct traffic classes: High-Priority (HP) packets representing critical protection alerts and fault notifications, and Normal-Priority (NP) packets representing telemetry and state-estimation feeds. HP generation follows an unannounced bursty Poisson process; NP traffic is periodic. Each node maintains separate FIFO buffers (HP capacity: 50 packets, NP capacity: 100). Packets exceeding deadlines ($T_H^{\text{max}} = 1.0$\,s, $T_N^{\text{max}} = 3.0$\,s) are dropped. The gateway scheduler operates at 100\,ms decision intervals over 5000-step episodes (500\,s real time).

\subsection{MDP Definition}
Scheduling is formulated as a discounted infinite-horizon MDP.
\begin{equation}
\mathbf{s}_t = \left[ q_t^{N},\; q_t^{H},\; e_t,\; a_t^{N},\; a_t^{H},\; d_t \right],
\end{equation}
where $q_t^{N}, q_t^{H} \in [0,1]$ denote normalized buffer occupancy; $e_t$ denotes normalized residual energy; $a_t^{N}, a_t^{H} \in [0,1]$ denote normalized head-of-line packet ages relative to expiration; and $d_t$ denotes normalized channel contention.

The action space comprises five discrete choices:
\begin{equation}
\mathcal{A} = \{\text{send HP (0)},\; \text{send normal (1)},\; \text{drop normal (2)},\; \text{sleep (3)},\; \text{drop HP (4)}\}.
\end{equation}
Actions 0, 1 attempt transmission; Action 3 enters sleep; Actions 2, 4 enable buffer purging of aged packets.

\subsection{Reward Formulation with Soft Policy Guidance}
Reward design addresses a fundamental challenge: naive throughput-based rewards induce policy collapse (ignoring critical traffic), while rigid penalties force replication of Strict Priority. The solution is outcome-driven reward with state-dependent guidance:
\begin{equation}
r_t = \tanh\left(\frac{r_t^{\text{base}} + r_t^{\text{guide}}}{100}\right),
\end{equation}
where base reward incorporates multiple objectives (confirmed deliveries with class-specific weights $w_1=25.0$ for HP, $w_2=5.0$ for NP; penalties for drops and timeouts; latency penalties; energy cost). The guidance term applies positive reinforcement for HP service and negative feedback for delayed HP packets only when urgent ($a_t^H > \tau_u = 0.55$), granting autonomy to trade brief critical delays for normal-queue recovery.
The base component reflects primary networking performance:
\begin{equation}
r_t^{\text{base}} = w_1 N_{\text{sent}}^H + w_2 N_{\text{sent}}^N + w_3 N_{\text{drop}}^H + w_4 N_{\text{drop}}^N + w_5 N_{\text{to}}^H + w_6 N_{\text{to}}^N + w_7 (a_t^H)^2 + w_8 (a_t^N)^2 - w_9 \Delta e_t,
\end{equation}
where $N_{\text{sent}}^H$ and $N_{\text{sent}}^N$ represent confirmed deliveries ($w_1 = 25.0, w_2 = 5.0$); $N_{\text{drop}}$ and $N_{\text{to}}$ penalize buffer overflow drops and timeout expirations ($w_3 = -40.0, w_4 = -8.0, w_5 = -80.0, w_6 = -10.0$); quadratic terms penalize waiting queue latencies ($w_7 = -40.0, w_8 = -2.0$); and $\Delta e_t$ penalizes energy depletion ($w_9 = 250.0$). If the scheduler selects sleep (Action 3), it receives a conservation bonus $r_{\text{sleep}} = +0.5$.

To ensure the agent reliably services critical traffic without eliminating operational flexibility, we introduce the soft policy guidance term:
\begin{equation}
r_t^{\text{guide}}=
\begin{cases}
+\alpha, & a_t = 0 \text{ and } q_t^H > 0\\
-\beta, & a_t \neq 0,\; q_t^H > 0,\; a_t^H > \tau_u\\
0, & \text{otherwise}
\end{cases}
\end{equation}
with bonus $\alpha = 6.0$, penalty $\beta = 8.0$, and urgency age threshold $\tau_u = 0.55$. Under this formulation, the agent receives positive reinforcement for servicing HP packets. Critically, the omission penalty $\beta$ is withheld if the waiting HP packet is young ($a_t^H \le \tau_u$). This design grants the scheduler autonomy to temporarily withhold a newly arrived HP packet to flush an aging normal queue, intervening with negative feedback only when an HP packet approaches critical delay boundaries.

\section{Methodology}

\subsection{Learning Architecture}
We implement Double Deep Q-Network (Double-DQN) to mitigate maximization bias in standard Q-learning. The target calculation separates action selection (via online network $\theta_t$) from evaluation (target network $\theta_t^{-}$):
\begin{equation}
Y_t^{\text{DoubleQ}} = r_{t+1} + \gamma Q\left(\mathbf{s}_{t+1}, \operatorname{arg\,max}_{a} Q(\mathbf{s}_{t+1}, a; \theta_t);\, \theta_t^{-}\right).
\end{equation}

Prioritized Experience Replay (PER) accelerates convergence on rare critical events, sampling transitions with probability $P(i) = p_i^\mu / \sum_k p_k^\mu$ where $p_i = |\delta_i| + \epsilon_{\text{per}}$ (temporal difference error) and $\mu = 0.6$. Importance weights $w_i = (N \cdot P(i))^{-\nu}$ correct for biased sampling, with $\nu$ annealing from 0.4 to 1.0.

The network is a fully connected MLP: input layer (6 neurons), two hidden layers (64 neurons, ReLU), output layer (5 Q-values). Target weights update via soft Polyak averaging with $\tau = 0.005$:
\begin{equation}
\theta^{-} \leftarrow \tau \theta + (1 - \tau)\theta^{-}.
\end{equation}
Optimization uses Adam ($\text{lr} = 1 \times 10^{-4}$, Cosine Annealing), gradient clipping (max $L_2 = 1.0$), and $\epsilon$-greedy exploration ($1.0 \to 0.05$ decay). Evaluation uses $\epsilon = 0.0$.

\subsection{Baseline Schedulers}
Comparisons employ two baselines: (i) \textbf{Strict Priority (SP)}: transmits HP when available, else NP; (ii) \textbf{Random}: uniformly samples actions.

\subsection{Training and Validation Protocol}
Training runs 50 epochs of 5000 steps. Checkpoint selection uses validation metric across fixed seeds:
\begin{equation}
S_{\text{val}} = 0.40 \left(\frac{\text{PDR}}{100}\right) + 0.10 \left(1 - \frac{\text{PLR}}{100}\right) + 0.35 \left(\frac{\text{HR}}{100}\right) + 0.15 \left(1 - \min\left(\frac{D_H}{1.0}, 1.0\right)\right).
\end{equation}
Early stopping triggers after 15 epochs without improvement. Evaluation executes 20 synchronized runs per policy (seeds 12545--12564) for rigorous statistical comparison.

\begin{table}[t]
\centering
\caption{Simulation and Training Hyperparameters (V76 Configuration).}
\label{tab:setup}
\scriptsize
\setlength{\tabcolsep}{6pt}
\begin{tabular}{ll}
\toprule
\textbf{Parameter Category} & \textbf{Configured Value} \\
\midrule
Network Simulator & ns-3.40 interfaced via ns3-gym \\
Network Topology & 100 Sensor Nodes + 1 Central Sink (150\,m $\times$ 150\,m) \\
Physical / MAC Layer & IEEE 802.15.4 (2.4\,GHz ISM Band, CSMA/CA) \\
Control Interval ($\Delta t$) & 100\,ms per decision slot \\
Episode Duration & 5000 decision steps (500\,s real time) \\
Traffic Generation & Bursty Poisson (HP Alerts) + Periodic Telemetry (NP) \\
Queue Capacities & HP Buffer: 50 packets; Normal Buffer: 100 packets \\
State / Action Dimensionality & $|\mathcal{S}| = 6$ continuous, $|\mathcal{A}| = 5$ discrete \\
Neural Architecture & MLP: $6 \to 64 \to 64 \to 5$ (ReLU activations) \\
Replay Buffer Capacity & 100,000 transitions (Prioritized: $\mu=0.6, \nu=0.4 \to 1.0$) \\
Mini-batch Size & 128 transitions \\
Discount Factor ($\gamma$) & 0.99 \\
Optimizer / Learning Rate & Adam ($\text{lr} = 1 \times 10^{-4}$, Cosine Annealing) \\
Target Network Update & Soft Polyak update ($\tau = 0.005$) \\
Exploration Schedule & $\epsilon$-greedy: $1.0 \to 0.05$ (decay factor $0.99995$) \\
Evaluation Repetitions & 20 matched runs per policy (seeds $12545$--$12564$) \\
\bottomrule
\end{tabular}
\end{table}

\subsection{Statistical Evaluation Protocol}
Metrics include aggregate throughput (kbps), Packet Delivery Ratio (PDR \%), Packet Loss Ratio (PLR \%), average latencies, energy per bit, and critical-packet delivery ratio. All are reported as mean ± 95\% confidence interval (Student's $t$). 

For pairwise significance testing, we verify normality of paired differences via Shapiro-Wilk test ($\alpha = 0.05$). If normal, paired $t$-test is applied; otherwise Wilcoxon signed-rank test. Effect sizes are quantified using Cohen's $d_z = \bar{D} / \sigma_D$ across the 20 paired runs.

\begin{table}[t]
\centering
\caption{WSN Performance Metrics Across 20 Matched Evaluation Runs (Mean $\pm$ 95\% CI).}
\label{tab:main_results}
\resizebox{\textwidth}{!}{%
\begin{tabular}{lccccccc}
\toprule
\textbf{Policy} & \textbf{Throughput} & \textbf{PDR} & \textbf{PLR} & \textbf{Normal Delay} & \textbf{HP Delay} & \textbf{Energy} & \textbf{HP Delivery} \\
 & \textbf{(kbps)} & \textbf{(\%)} & \textbf{(\%)} & \textbf{(s)} & \textbf{(s)} & \textbf{(nJ/bit)} & \textbf{Ratio (\%)} \\
\midrule
Proposed DQN-Edge & $7.0267 \pm 0.0702$ & $50.3291 \pm 0.6022$ & $40.3309 \pm 0.8373$ & \textbf{0.4287 $\pm$ 0.0261} & $0.1431 \pm 0.0026$ & $443.8708 \pm 16.9422$ & $95.7779 \pm 0.2916$ \\
Strict Priority (SP) & \textbf{7.9104 $\pm$ 0.0054} & \textbf{56.6625 $\pm$ 0.3137} & \textbf{34.7224 $\pm$ 0.5326} & $0.6171 \pm 0.0384$ & \textbf{0.0274 $\pm$ 0.0015} & \textbf{417.0370 $\pm$ 17.2341} & \textbf{98.3374 $\pm$ 0.1941} \\
Random Baseline & $2.8000 \pm 0.0350$ & $20.0709 \pm 0.3314$ & $71.4943 \pm 0.5450$ & $0.8545 \pm 0.0597$ & $0.3521 \pm 0.0129$ & $806.7023 \pm 22.7489$ & $34.4862 \pm 0.8563$ \\
\bottomrule
\end{tabular}%
}
\end{table}

\begin{table}[ht]
\centering
\caption{Matched Pairwise Statistical Comparison: Proposed DQN-Edge vs. Strict Priority ($n=20$ runs).}
\label{tab:stats}
\scriptsize
\setlength{\tabcolsep}{8pt}
\begin{tabular}{llcc}
\toprule
\textbf{Evaluated Metric} & \textbf{Statistical Test} & \textbf{\textit{p}-value} & \textbf{Paired Cohen's $d_z$} \\
\midrule
Throughput (kbps) & Paired $t$-test & $2.656 \times 10^{-16}$ & $-5.8046$ \\
PDR (\%) & Paired $t$-test & $1.655 \times 10^{-15}$ & $-5.2568$ \\
PLR (\%) & Paired $t$-test & $4.749 \times 10^{-13}$ & $+3.8504$ \\
Avg. Normal Delay (s) & Wilcoxon signed-rank & $1.907 \times 10^{-6}$ & $-3.3284$ \\
Avg. HP Delay (s) & Paired $t$-test & $3.038 \times 10^{-26}$ & $+19.6001$ \\
Energy (nJ/bit) & Paired $t$-test & $1.687 \times 10^{-14}$ & $+4.6310$ \\
HP Delivery Ratio (\%) & Paired $t$-test & $2.176 \times 10^{-11}$ & $-3.1005$ \\
\bottomrule
\end{tabular}
\end{table}

\subsection{Convergence and Training Stability}
Figure~\ref{fig:learning} illustrates the training progression over 30 elapsed epochs prior to early stopping. The validation score trajectory shows rapid initial policy improvement over the first 5 epochs, transitioning into stable, monotonic refinement without policy collapse or loss divergence. The exploration rate $\epsilon$ decays steadily, confirming that the agent transitions from exploratory state-space coverage to stable deterministic execution. The model achieving the highest validation score at Epoch 12 was preserved for full benchmarking.

\begin{figure}[htbp]
    \centering
    \includegraphics[width=0.72\columnwidth]{figures/fig_1_learning_curve.png}
    \vspace{-2mm}
    \caption{Double-DQN training progression and validation-score trend over 30 epochs (V76).}
    \label{fig:learning}
\end{figure}

\subsection{Comparative Analysis: Understanding the Trade-Off}
Against Random, our DQN agent achieves 150.9\% higher throughput and 95.78\% critical delivery (vs. 34.49\%), with 45\% lower energy waste, confirming multi-class learning success. Strict Priority comparison shows different characteristics. Strict Priority achieves superior performance on latency-critical metrics: 7.9104\,kbps versus 7.0267\,kbps throughput, 0.0274\,s versus 0.1431\,s critical latency (5.2$\times$ faster), and 98.34\% versus 95.78\% delivery ratio (Table~\ref{tab:stats}, $p < 10^{-11}$). The statistical margin is substantial and unambiguous.

The learned policy exhibits different behavior when grid disturbances sustain high alert volumes. Strict Priority exhausts transmission capacity on urgent traffic, starving normal telemetry. During these episodes, operators lose access to state-estimation feeds and monitoring exactly when awareness is most needed. The learned scheduler recovers 30.5\% of lost telemetry responsiveness (0.6171\,s $\rightarrow$ 0.4287\,s) by accepting 2.56\% lower critical-delivery performance.

\subsection{Delay Distribution and Buffer Dynamics}
Figures~\ref{fig:bar} and~\ref{fig:cdf} present the mean performance metrics with 95\% confidence intervals and the empirical Cumulative Distribution Functions (CDFs) of packet latencies for both traffic classes.

\begin{figure}[htbp]
    \centering
    \begin{minipage}[b]{0.48\textwidth}
        \centering
        \includegraphics[width=\textwidth]{figures/fig_4_bar_charts.png}
        \caption{Performance metrics (Mean $\pm$ 95\% CI).}
        \label{fig:bar}
    \end{minipage}
    \hfill
    \begin{minipage}[b]{0.48\textwidth}
        \centering
        \includegraphics[width=\textwidth]{figures/fig_5_delay_cdf.png}
        \caption{Delay CDFs for traffic classes.}
        \label{fig:cdf}
    \end{minipage}
    \vspace{-2mm}
\end{figure}

Analysis of the empirical delay CDFs in Figure~\ref{fig:cdf} illuminates the underlying mechanics of each scheduling discipline. Strict Priority aggressively prioritizes high-priority service, clearing over 90\% of all critical packets within 50\,ms. However, this aggressive service strategy induces severe queue starvation on normal traffic, manifesting as a long-tailed delay distribution averaging 0.6171\,s. The learned policy accepts higher critical latency (0.1431\,s mean) to shift normal-traffic delays leftward in the distribution. Table~\ref{tab:main_results} quantifies this trade-off: the learned policy achieves 30.5\% lower average normal delay than Strict Priority ($p = 1.907 \times 10^{-6}, d_z = -3.3284$).

\subsection{Trade-Off Analysis in Smart Grid Telemetry}
The empirical results demonstrate a clear engineering trade-off between conflicting objectives. Strict Priority achieves ultra-low critical latency (0.0274\,s) at the cost of severe normal-queue starvation (0.6171\,s). The learned DQN policy operates at a different point on this trade-off frontier. Through its state representation (particularly the head-of-line packet ages $a_t^H$ and $a_t^N$), the agent identifies windows when newly arrived critical packets retain substantial deadline margin. During these windows, the scheduler opportunistically serves aging normal packets before critical deadlines force urgent service. This learned strategy exchanges a 2.56\% reduction in critical-packet delivery for a substantial 30.5\% improvement in normal-traffic latency. For grid operators managing sustained disturbances, this fairness improvement directly supports situational awareness maintenance.

\subsection{Policy Collapse and Reward Shaping}
During development, early configurations (V74) exhibited catastrophic policy collapse: the learned policy transmitted normal packets at 100\% frequency while delivering zero critical packets. Initial investigation attributed this to training dynamics. Closer examination revealed that the simulator's safeguard against zero-division errors (defaulting average delays to 0.0 when no packets delivered) masked this degenerate behavior during high-throughput phases. The policy exploited this implementation detail, achieving superficially high rewards while failing completely on the critical traffic objective. Transitioning to explicit per-class validation metrics in V76 resolved the issue. Separating class-specific metrics and applying penalties only when HP packets exceed urgency thresholds ($a_t^H > 0.55$) prevents policy collapse while maintaining acceptable dual-objective performance. This finding highlights an important constraint: mixed-criticality RL requires explicit per-class validation and tail-latency visibility. Aggregate reward signals miss critical failure modes like complete traffic-class starvation.

\subsection{Sensitivity Analysis: Cross-Configuration Robustness}
In the 50-node scenario (reduced congestion), the trained DQN achieved 0.3812\,s normal delay versus Strict Priority's 0.5421\,s (29.7\% improvement). At 200 nodes (elevated congestion), the respective values were 0.5193\,s versus 0.7846\,s (33.8\% improvement). The consistent improvement across topologies indicates that the learned fairness principle transfers to different network scales, although sensitivity to traffic ratios, deadline configurations, and deployment specifics requires further investigation.

\section{Discussion}
Double-DQN produces adaptive scheduling with multiclass fairness properties. However, Strict Priority achieves 5.2× lower critical latency (0.0274\,s vs. 0.1431\,s), exposing inherent limitations of the learning approach: the compact state representation constrains observability of fine-grained queue dynamics, and single-agent Q-learning cannot match the deterministic responsiveness of rule-based heuristics optimized specifically for latency minimization. Bridging this gap requires explicit safety-layer constraints (e.g., Constrained MDPs) and distributional value functions for tail-latency control, neither of which were incorporated in this study.

\textbf{Baseline Scope:} This work compares learned scheduling against Strict Priority and Random baselines. The absence of adaptive rule-based comparisons (Weighted Round-Robin, Earliest Deadline First, Weighted Fair Queuing) limits the scope of conclusions regarding where learning-based methods specifically outperform classical heuristics. Such comparisons constitute an important avenue for future investigation to clarify the competitive advantage of DRL in this domain.

\textbf{Advancing Safe DRL in Comms:} A key observation from this analysis: learned policies do not universally outperform rule-based methods in safety-critical networks. This suggests an important requirement for further work: learning-based methods require explicit constraints and formal verification prior to deployment. Promising technical directions include Constrained MDPs for safety layer enforcement, distributional value functions (QR-DQN, IQN) for tail-latency control, and formal verification frameworks to provide mathematical guarantees on worst-case performance.

\section{Threats to Validity}
All evaluation took place within discrete-event simulation (ns-3.40). While ns-3 provides high-fidelity 802.15.4 CSMA/CA modeling, physical deployments introduce non-idealities: jitter, clock skew, multipath fading, asymmetric propagation. This simulation-to-reality gap could degrade algorithm performance. Second, the observed trade-off is specific to the tested scenario (100 nodes, 100\,ms intervals, fixed deadlines). Varying traffic proportions or node density redistributes the Strict Priority versus learning performance frontier. Finally, the approach addresses single-hop cluster scheduling. Multi-hop mesh topologies introduce distributed coordination requirements not considered here.

\section{Conclusion}
This work contributes a reproducible empirical analysis of policy-guided Double-DQN for mixed-criticality packet scheduling in smart grid sensor networks. Through systematic comparison across 20 matched evaluation trials using ns-3 simulation coupled with PyTorch, we establish a quantitative trade-off: the learned scheduler sacrifices 2.56\% critical-packet delivery and 115.7\,ms critical latency to recover 30.5\% of the normal-traffic responsiveness that rule-based Strict Priority scheduling loses. This fairness improvement supports operator situational awareness during grid disturbances when sustained telemetry visibility directly impacts emergency response quality. Benchmarking establishes where learning-based methods add value and where rule-based systems show advantages. This analysis sets expectations for deploying RL in safety-critical settings and identifies needed technical contributors (explicit constraints, distributional value functions, formal verification).

\section*{Reproducibility Package}
To facilitate independent verification, the complete simulation artifacts---including the raw evaluation results (\texttt{raw\_results\_v76.csv}), training history (\texttt{training\_log\_v76.csv}), statistical summary tables (\texttt{final\_results\_v76.xlsx}), and the underlying ns-3 C++ and PyTorch source code---are made available for academic reuse.

\bibliographystyle{splncs04}
\bibliography{references}

\end{document}